\section{The two fine theories in categories and simplicial sets}

\begin{claim}{B.1}{Proved}{The comma tests and their variance}
For a fundamental localizer $W$ and a small category $S$, put
\[
\begin{aligned}
 W^{\mathrm{pre}}_S&=\{u:X\to Y\text{ over }S:\ s\backslash u\in W\text{ for all }s\in S\},\\
 W^{\mathrm{co}}_S&=\{u:X\to Y\text{ over }S:\ u/s\in W\text{ for all }s\in S\}.
\end{aligned}
\]
Both are weakly saturated classes contained in the total class $W_S$.
Opposition gives an isomorphism of relative categories
$(\Cat/S,W^{\mathrm{pre}}_S)\cong
(\Cat/S^{\op},W^{\mathrm{co}}_{S^{\op}})$.
For $W=W_\infty$, their enhanced localizations are, respectively,
$\Fun(NS^{\op},\Spaces)$ and $\Fun(NS,\Spaces)$.
The value of the first classifier at $s$ is $B(s\backslash X)$;
the value of the second is $B(X/s)$.
In the notation of XII ms.~222, $W^{\mathrm g}_S$ is the second class.
\end{claim}
\begin{proof}
Each comma construction is a functor on $\Cat/S$, so identities, two-out-of-three,
and retracts are inherited from $W$. Relative Theorem A gives
$W^{\mathrm{co}}_S\subset W_S$, and its opposite form gives the other inclusion.
There are canonical category isomorphisms
$(s\backslash X)^{\op}\cong X^{\op}/s$; fundamental localizers are invariant
under opposition \citep[Sorites 3.3.8]{BCCisinski2006}. These are isomorphisms
of specified comma categories, not just equivalences of their nerves.
The enhanced identifications are proved in B.3 below. An arrow $s\to t$
acts by precomposition $t\backslash X\to s\backslash X$ and by
postcomposition $X/s\to X/t$, which fixes the variance. For $X=\varnothing$
all values are empty; for $S=\varnothing$ the slice has its unique object.
\end{proof}

\begin{claim}{B.2}{Known}{The simplicial slice test, with the necessary base hypothesis}
Let $S$ be a small quasicategory and let $u:X\to Y$ be any map of simplicial
sets over $S$; neither $X$ nor $Y$ is assumed fibrant. Then $u$ is a covariant
equivalence exactly when
\[
 X\times_S S_{/s}\longrightarrow Y\times_S S_{/s}
\]
is a weak homotopy equivalence for every vertex $s$. Replacing $S_{/s}$ by
$S_{s/}$ gives the contravariant test. For an arbitrary simplicial base these
statements hold with a right, respectively left, fibrant replacement
$\widetilde S_{/s}\to S$, respectively $\widetilde S_{s/}\to S$, of the
vertex map in place of the literal slice.
\end{claim}
\begin{proof}
The precise all-base statement is \citet[4.4.29--4.4.30]{BCCisinski2019}.
We recall why the replacement and the word ``strict'' matter. Choose a left
anodyne map $X\to X^{\mathrm{lf}}$ over $S$ with left-fibration target.
Pulling it back along the right fibration $\widetilde S_{/s}\to S$ remains
a covariant equivalence, hence a total weak homotopy equivalence. The same
holds for $Y$. The inclusion of the fibre
$(X^{\mathrm{lf}})_s\to X^{\mathrm{lf}}\times_S\widetilde S_{/s}$ is a
weak equivalence: the chosen vertex is final in $\widetilde S_{/s}$ and a
left fibration preserves final maps under pullback. Consequently the displayed
test is precisely the fibre test for a map of left fibrations. That test detects
equivalences by \citet[Theorem 4.4.14]{BCCisinski2019}, or
\citet[Theorem 2.2.1.2 and Corollary 2.2.3.12]{DGLurieHTT}.
When $S$ is a quasicategory, $S_{/s}\to S$ itself is such a right-fibrant
replacement, with the vertex $1_s$ final. This supplies the actual comparison
maps used in the reduction. Opposition proves the other assertion.
\end{proof}

\begin{claim}{B.3}{Proved}{Nerve comparison as an equivalence of enhanced localizations}
For every small ordinary category $S$, the ordinary nerve induces equivalences
\[
 (\Cat/S)[(W^{\mathrm{co}}_S)^{-1}]_\infty
 \simeq (\sSet/NS)[(W^{\mathrm{cov}})^{-1}]_\infty
 \simeq\Fun(NS,\Spaces),
\]
and, by opposition,
\[
 (\Cat/S)[(W^{\mathrm{pre}}_S)^{-1}]_\infty
 \simeq (\sSet/NS)[(W^{\mathrm{contra}})^{-1}]_\infty
 \simeq\Fun(NS^{\op},\Spaces).
\]
Here $W=W_\infty$. These are equivalences of infinity-categories and therefore
of homotopy categories. The last comparison is implemented by the
straightening--unstraightening Quillen equivalence. No model structure on
$\Cat/S$ with the ordinary nerve as a right Quillen functor is asserted.
\end{claim}
\begin{proof}
The nerve preserves strict pullbacks and has
$N(S/s)=(NS)_{/s}$ and $N(s\backslash S)=(NS)_{s/}$. Thus B.2 identifies
exactly the weak equivalences on nerves. We verify essential surjectivity and
all coherent mapping information by giving inverse relative functors.
For $K\to NS$, let $T(K)=\Delta/K$ and map it to $S$ by the last vertex.
Its objects are simplices $\sigma:\Delta^n\to K$; a morphism
$\sigma\to\tau$ is a monotone map $\theta:[n]\to[m]$ with
$\sigma=\tau\theta$. Its image in $S$ is the last segment
$\tau(\theta(n))\to\tau(m)$. Composition is composition of monotone maps.
There is a natural last-vertex map $\epsilon_K:NT(K)\to K$ over $NS$.
For each $s$ the strict identification
\[
 T(K)/s\cong\Delta/(K\times_{NS}N(S/s))
\]
sends $(\sigma,\alpha:\sigma(n)\to s)$ to the simplex of $S/s$ whose last
arrow is $\alpha$. In the other direction a simplex of $S/s$ is determined
by its last such arrow. This also checks the morphisms.
Under this identification the pullback of $\epsilon_K$ is the usual
last-vertex map from the nerve of a simplex category. That map is a weak
homotopy equivalence: the simplicial replacement of the simplex decomposition
has contractible simplex fibres, and its two augmentations to
$N(\Delta/L)$ and $L$ are weak equivalences. Equivalently, this is the
simplex-category instance of \citet[Proposition 7.3.8 and Examples 7.3.9--7.3.10]{BCCisinski2019}.
It follows from B.2 that every $\epsilon_K$ is a covariant equivalence.
Two-out-of-three now proves that $T$ preserves covariant equivalences into
$W^{\mathrm{co}}_S$.

For a category $X\to S$, evaluation at the last vertex is a functor
$T(NX)\to X$ over $S$, and the preceding calculation shows that it is in
$W^{\mathrm{co}}_S$. The two natural transformations
$NT\Rightarrow1$ and $TN\Rightarrow1$ become natural equivalences after
infinity-localization. Thus the two localized functors are inverse, including
their mapping spaces. Opposition gives the presheaf comparison. Finally,
straightening over the ordinary nerve identifies the covariant localization
with coherent $S$-diagrams of spaces; its Quillen equivalence is
\citet[Theorem 2.2.1.2]{DGLurieHTT}. For the first classifier the same
argument, or the pointwise left Kan formula applied to the Yoneda generators,
gives exactly $s\mapsto B(s\backslash X)$.
\end{proof}

\begin{claim}{B.4}{Refuted}{The literal slice recipe over every simplicial base}
Lead~1 is false for arbitrary simplicial bases when its tests use literal
simplicial slices. Let $S=\Lambda^2_1$, with nondegenerate edges $01$ and
$12$ and no $02$ edge. The map $\{0\}\hookrightarrow\Delta^{\{0,1\}}$
over $S$ is a covariant equivalence, but its pullback along $S_{/2}\to S$
is $\varnothing\to\{1\}$ and is not a weak homotopy equivalence.
\end{claim}
\begin{proof}
The given map is the generating left-anodyne inclusion
$\Lambda^1_0\hookrightarrow\Delta^1$, hence is a covariant equivalence
over every base \citep[Proposition 2.1.4.9]{DGLurieHTT}.
The vertices of $S_{/2}$ are the edge $12$ and the degenerate edge $22$;
their projections are $1$ and $2$. There is an edge between them supplied by
the degeneracy of $12$, but none projecting to $01$. Therefore its pullback
over $\{0\}$ is empty and over the edge $01$ is the vertex $1$.
All higher simplices in the latter are degenerate. The two numbers of
components are zero and one. The finite simplex data and both intersections
are checked in \texttt{code/bc\_counterexamples.py}.
Applying simplicial opposition gives the contravariant counterexample.
This refutes the assertion as stated, while B.2 supplies the corrected version.
\end{proof}

\begin{claim}{B.5}{Known}{The exact dictionary for smoothness and properness}
For $p:A\to B$ in $\Cat$, Grothendieck's $W_\infty$-properness is equivalent
to properness of $Np$ in \citet[Definition 4.4.1]{BCCisinski2019}, and to
smoothness of $Np$ in \citet[Definition 4.1.2.9]{DGLurieHTT}.
Grothendieck's $W_\infty$-smoothness corresponds to Cisinski-smoothness and
to the opposite of Lurie's smoothness. Thus the two cited books use opposite
names. The equivalences hold in both directions for all ordinary functors,
without a fibration hypothesis.
\end{claim}
\begin{proof}
For an arrow $\alpha:p(a)\to b$, define the outgoing lift category
$L^+_p(a,\alpha)$: an object is $v:a\to a'$ with $p(v)=\alpha$; an arrow
$v\to v'$ is $w:a'\to a''$ with $p(w)=1_b$ and $wv=v'$.
It is the comma category $(a,\alpha)\backslash A_b$ for
$A_b\to A/b$. Thus Grothendieck-properness says precisely that all these
categories have weakly contractible nerves. The incoming version is dual.
Every nerve of an ordinary functor is an inner fibration, since the unique
inner-horn fillers in nerves are the specified composites. In the ordinary
nerve case, \citet[Theorem 4.4.36, condition (ii')]{BCCisinski2019}
identifies the displayed condition with simplicial properness. This particular
case is unchanged in the author's July 2026 errata
\citep[pp.~1--4]{BCCisinskiErrata2026}; for arbitrary quasicategories the
corresponding fibre-to-comma condition must also be imposed after all base
changes. We do not use the incorrect unrestricted version of that theorem.
Finally Cisinski-properness is universal preservation of final maps,
equivalently of right-anodyne maps (Proposition~4.4.4). HTT
Proposition~4.1.2.8 makes this exactly Lurie-smoothness. Opposition supplies
the second equivalence. In particular $\{1\}\to[1]$ is
Grothendieck-proper and Lurie-smooth, whereas $\{0\}\to[1]$ has the other
two names. Neither class is the class of proper maps of geometric realizations.
\end{proof}

\begin{claim}{B.6}{Proved}{The maximal localization class, for every fundamental localizer}
Let $\widehat{\mathcal L}_{\mathrm{co}}$ consist of those $p:A\to B$ for
which every strict base change $p_T:A\times_BT\to T$ preserves
$W^{\mathrm{co}}_T$ under raw pullback. Then
\[
 \widehat{\mathcal L}_{\mathrm{co}}=\Liss_W,
 \qquad \widehat{\mathcal L}_{\mathrm{pre}}=\Prop_W.
\]
No accessibility assumption is needed. In particular the equality left open
at XII ms.~222 has an affirmative answer in the manuscript's covariant
variance; it is not the assertion that all raw inverse images are homotopical.
\end{claim}
\begin{proof}
We use the following exact form of \citet[Theorem 1.2(c)]{BCMaltsiniotis2005}:
a functor is $W$-smooth if and only if all its base changes preserve
$W$-aspherical functors under pullback. Here ``aspherical functor'' means
that all its over-comma categories are $W$-aspherical.
For $X\to T$, let $R_TX=\int_{t\in T}(X/t)$, with its opfibration to $T$,
and let $j_X:X\to R_TX$ send $x$ to $(p(x),x,1_{p(x)})$.
An object of $R_TX$ is $(t,x,\beta:p(x)\to t)$. The category
$j_X/(t,x,\beta)$ is canonically $X/x$: an arrow from $j_X(y)$ to the
given object is determined by $v:y\to x$, since its base arrow must be
$\beta p(v)$. Thus $j_X$ is $W$-aspherical for every fundamental localizer.
If $q:T'\to T$ is smooth, its pullback
$(R_TX)\times_TT'\to R_TX$ is smooth. The cited characterization implies
that $j'_X:X\times_TT'\to(R_TX)\times_TT'$ is aspherical.
An aspherical functor $v:U\to V$ over $T'$ is in
$W^{\mathrm{co}}_{T'}$: for each $t'$ and $(z,\gamma)\in V/t'$, the
comma category $(v/t')/(z,\gamma)$ is $v/z$, so relative Theorem A first
makes $v/t'$ aspherical and then a $W$-map.
For $u:X\to Y$ in $W^{\mathrm{co}}_T$, the map $R_Tu$ is fibrewise in $W$.
Its pullback along $q$ is still a fibrewise $W$-map of opfibrations.
Over-comma categories of an opfibration retract, by cocartesian transport,
onto its fibres: the inclusion of the fibre over $t'$ has a left adjoint
sending $(z,\alpha:pz\to t')$ to $\alpha_!z$. Hence this pulled-back map
lies in $W^{\mathrm{co}}_{T'}$.
In the natural square with vertical arrows $j'_X,j'_Y$, three sides are now
in $W^{\mathrm{co}}_{T'}$; two-out-of-three gives $q^*u$ in that class.
Smoothness is preserved by arbitrary base change by its lifting criterion,
which proves $\Liss_W\subset\widehat{\mathcal L}_{\mathrm{co}}$.

Conversely, a functor $u:X\to T$ is $W$-aspherical precisely when the map
$(X\to T)\to(T\xrightarrow{1}T)$ is in $W^{\mathrm{co}}_T$, since
$T/t$ has a terminal object. If $p$ belongs to the universal-preservation
class, every pullback of every such aspherical functor is aspherical.
This is exactly condition (c) of the same Theorem~1.2, hence $p$ is smooth.
This argument also proves maximality, not merely containment. Opposition
proves the presheaf assertion. It explains why interchanging $s\backslash X$
and $X/s$ while leaving the name ``smooth'' unchanged gives a false claim.
\end{proof}

\begin{claim}{B.7}{Proved}{Three different fibration classes at the minimal localizer}
There are strict inclusions
\[
 \{\text{Grothendieck bifibrations}\}\subsetneq
 \Liss_{W_\infty}\cap\Prop_{W_\infty}\subsetneq\Fib_{W_\infty}.
\]
For simplicial sets, Cisinski-smooth and Cisinski-proper maps likewise lie in
the intrinsic class preserving all weak homotopy equivalences under strict
pullback. The analogous inclusion for a larger fundamental localizer is
generally false.
\end{claim}
\begin{proof}
A cartesian lift supplies a terminal object in each incoming lift category;
a cocartesian lift supplies an initial object in each outgoing lift category.
This proves the first inclusion. Cisinski's Corollary~4.6.5 and
Proposition~4.6.2 identify a simplicial map which is both smooth and proper
as one whose every strict pullback square is a homotopy pullback. Homotopy
pullbacks preserve weak equivalences. B.5 and preservation of pullbacks by
the nerve prove the second inclusion for categories.

For strictness of the first inclusion use the four-object poset
$x<a<c$, $x<b<c$, with $a,b$ incomparable, mapped to $[1]$ by
$x\mapsto0$ and $a,b,c\mapsto1$. The only nontrivial outgoing lift poset is
$\{a,b<c\}$, which has terminal object $c$ but no initial object. All incoming
nonidentity tests are singletons. It is therefore smooth and proper but has
no cocartesian lift from $x$ over $0\to1$.
For strictness of the second inclusion let $V$ be the three-object poset
$a<b$, $a<c$, and map $a,b$ to $0$ and $c$ to $1$.
The section $s$ given by $s(0)=a$, $s(1)=c$ satisfies $ps=1$ and
there is a natural transformation $sp\Rightarrow1_V$ over $[1]$.
It survives every strict pullback. Thus each pullback of $p$ is a weak
equivalence, and two-out-of-three in the natural square with any weak
equivalence over $[1]$ proves $p\in\Fib_{W_\infty}$.
There is no outgoing lift from $b$ to colour $1$, so $p$ is not proper.
The program checks every finite category law, all diamond lift posets,
and the natural-retraction data. These examples also prove the simplicial
strictness statements after taking nerves.
\end{proof}

\begin{claim}{B.8}{Proved}{An explicit smooth and proper path resolution}
For a small category $C$ and $x\in C$, let $E_x(C)$ have as objects the
functors $\gamma:J\to C$, where $J$ has objects $\mathbb Z$ and arrows
$2i\to2i-1$, $2i\to2i+1$, which are constantly $x$ with identity arrows
at the left end and constantly some object at the right end. Morphisms are
natural transformations whose eventual left component is $1_x$; composition
is componentwise. Then $q:E_x(C)\to C$, taking the eventual right value,
is $W_\infty$-smooth and proper, $BE_x(C)$ is weakly contractible, and its
strict fibres are the homotopy fibres of $\{x\}\to BC$.
\end{claim}
\begin{proof}
All objects and morphisms are small: eventually constant diagrams form a
small union of finite-diagram sets. Natural transformations have an eventual
right component because every sufficiently remote horizontal arrow is an
identity. Thus $q$ is an ordinary functor.
Fix $\gamma$ and an arrow $\alpha:q\gamma\to z$. For every sufficiently
large even integer $k$, let $\gamma_k$ agree with $\gamma$ through $k$ and
be constantly $z$ after $k$, with the arrow at $k\to k+1$ equal to $\alpha$.
There is a natural transformation $\eta_k:\gamma\to\gamma_k$, identity
through $k$ and equal to $\alpha$ afterwards. In the outgoing lift category,
restrict to those $(\eta:\gamma\to\delta)$ whose target is stationary after
$k+1$. This full subcategory has initial object $\eta_k$: the unique map
$\gamma_k\to\delta$ under $\gamma$ is $\eta_i$ for $i\leq k$ and
$1_z$ thereafter. Naturality at the joining edge is exactly the naturality
of $\eta$. The eventual identity condition forces uniqueness on the entire
tail. These subcategories, as $k$ increases through even integers, cover
every object and morphism of the outgoing lift category. Their nerves form
a filtered union of contractible simplicial sets, hence are weakly
contractible: every finite simplicial sphere and null-homotopy question occurs
at a finite stage. For an incoming arrow $z\to q\gamma$, choose $k$ odd,
use the joining arrow $k+1\to k$, and obtain terminal objects in the analogous
full subcategories. This proves smoothness and properness directly.

The entire category $E_x(C)$ is the filtered union of the categories of
finite-fence diagrams on $[-n,n]$ with leftmost value $x$, extended constantly
outside that interval. Each finite-fence category is contractible. Indeed,
forgetting its last vertex has constant extension as a left or right adjoint,
according to the orientation of the last edge; repeat until only the fixed
object remains. These are actual adjunctions, with components given by the
last edge. The same finite-stage argument proves $BE_x(C)\simeq *$.
B.7 now says every strict pullback along $q$ is a homotopy pullback.
The constant path gives the specified factorization of the chosen point and
identifies its fibres with the homotopy fibres of that point map.
The construction is an integer-fence version of the categorical cocylinder
\citep{BCEvrard1975}; the proof here does not assume the inherited PATHS claims.
\end{proof}

\begin{claim}{B.9}{Proved}{The exact localizer threshold for complementarity}
For a fundamental localizer $W$ in the relative-Theorem-A convention,
\[
 \Liss_W\cap\Prop_W\subset\Fib_W
 \quad\Longleftrightarrow\quad
 W=W_\infty\ \text{or}\ W=\operatorname{Mor}(\Cat).
\]
Accessibility is not required. Thus the tentative complementarity at
XII ms.~225--226 is false for every proper enlargement of $W_\infty$
other than the total class.
\end{claim}
\begin{proof}
Minimality gives $W_\infty\subset W$
\citep[Corollary 4.2.19]{BCCisinski2006}. The lift criteria show
$\Liss_{W_\infty}\subset\Liss_W$ and
$\Prop_{W_\infty}\subset\Prop_W$.
Every functor $v:C\to B$ has a factorization
$C\xrightarrow{i}P(v)\xrightarrow{p}B$ with $i\in W_\infty$ and $p$
smooth and proper for $W_\infty$: apply the verified model structure of
\citet[Theorem 5.3.14 and Proposition 5.3.13]{BCCisinski2006} only to the
accessible minimal localizer. By B.7 every strict pullback along $p$ is a
homotopy pullback after realization. If the displayed inclusion holds,
$p$ lies in $\Fib_W$. For $u:A\to B$ in $W$, the map
$A\times_BP(v)\to P(v)$ therefore lies in $W$ and realizes the homotopy
pullback of $Bu$ along $Bv$.

This proves that the class induced by $W$ in spaces is closed under arbitrary
homotopy pullback. Here the passage between models has specified comparisons:
any strict simplicial cospan is replaced by its simplex categories, whose
last-vertex comparison maps are weak equivalences; replacing one leg by $p$
then computes its homotopy pullback. Two-out-of-three transfers $W$ across
these comparisons. The fundamental-localizer/simplicial-localizer
correspondence \citep[Theorem 4.2.15]{BCCisinski2006} ensures this uses the
same class, also when it is not accessible.

Suppose $u\in W\setminus W_\infty$. Some homotopy fibre $F$ of $Bu$ is
not weakly contractible, and pullback stability makes $F\to*$ a $W$-map.
If $F=\varnothing$, we already have $\varnothing\to*\in W$.
Otherwise each point $*\to F$ is in $W$ by two-out-of-three; pulling it
back along itself shows $\Omega F\to*$ is in $W$. Repeating, a nonzero
homotopy group or an initial failure of connectedness gives a $W$-aspherical
space $Z$ with two components. The homotopy pullback of two points in
different components is empty, again yielding $\varnothing\to*\in W$.
Products with arbitrary categories then give $\varnothing\to A\in W$,
and two-out-of-three gives every $A\to B$ in $W$.
Conversely the total class satisfies the inclusion trivially and $W_\infty$
satisfies it by B.7. The rigidity agrees with
\citet[Theorem 6.4.22 and Remark 6.4.23]{BCCisinski2006}; the argument just
given supplies the exact no-accessibility and empty-category qualifications.
\end{proof}

\begin{claim}{B.10}{Refuted}{Complementarity for every finite truncation localizer}
For every $n\geq0$, $\Liss_{W_n}\cap\Prop_{W_n}$ is not contained in
$\Fib_{W_n}$. For $n=0$ a finite witness is the regular covering
$EC_2\to BC_2$. For $n\geq1$ B.8 supplies a completely specified witness
with base the face poset of the boundary of an $(n+2)$-simplex.
\end{claim}
\begin{proof}
The category $EC_2$ has objects $0,1$, a unique arrow between every ordered
pair, and forced composition. Send $i\to j$ to $j-i\pmod2$ in the one-object
group category $BC_2$. Every arrow has a unique outgoing and incoming lift,
so the map is a discrete bifibration and smooth and proper for every $W$.
The point map $e\to BC_2$ is in $W_0$, but its pullback has two discrete
objects and maps to the connected category $EC_2$. It is not in $W_0$.
All laws and component counts are checked by the finite program.
For general $n$, let $C_n$ have the nonempty proper subsets of
$\{0,\ldots,n+2\}$ as objects, inclusions as arrows, and inclusion composition.
Its nerve is the barycentric subdivision of $S^{n+1}$. Choose any vertex $x$.
The map $e\to C_n$ belongs to $W_n$. Pullback along $E_x(C_n)\to C_n$
has fibre of type $\Omega S^{n+1}$, and total category weakly contractible.
For $n\geq1$, its $\pi_n$ is $\mathbb Z$, so the resulting map is not in
$W_n$. Every object, arrow and composition of this infinite counterexample
was specified in B.8; the finite-base sphere assertion does not assert a
finite total category. B.12 gives a finite total-category witness for $n=1$.
\end{proof}

\begin{claim}{B.11}{Refuted}{A single covering refutes complementarity for every nonzero homology theory}
Let $E$ be any nonzero spectrum, with no connectivity assumption, and let
$W_E$ mean equivalence on $E\wedge\Sigma^\infty(-)_+$.
The three-sheeted action-groupoid covering in the prompt is both
$W_E$-smooth and $W_E$-proper, but is not in $\Fib_{W_E}$.
\end{claim}
\begin{proof}
Let $F=\langle x,y\rangle$ be the free group and let it act on
$\{1,2,3\}$ by $x=(12)$, $y=(13)$. The total category has these three
objects and an arrow labelled $g\in F$ from $s$ to $g(s)$; composition is
multiplication in $F$. Its projection to the one-object category $BF$ is
a discrete fibration and opfibration. Each lifting category is a singleton,
so both $W_E$ conditions hold for the same localizer.
Define $\alpha:F\to F$ by $\alpha(x)=x$ and
$\alpha(y)=yxyx^{-1}y^{-1}$. On the graph $BF\simeq S^1\vee S^1$, the
stable reduced suspension-spectrum map of $B\alpha$ has matrix $I_2$,
because the exponent sums are $(1,0)$ and $(0,1)$. Adjoining the basepoint
adds an identity sphere summand. Hence $B\alpha\in W_E$ for every spectrum.
The original action is transitive, whereas the pulled-back action sends
both generators to $(12)$ and has orbits $\{1,2\}$ and $\{3\}$.
The pullback map therefore identifies two components. To see the failure
for nonconnective $E$ as well, choose $k$ with $\pi_kE\ne0$. Point classes
in the two source components give a nonzero difference, split by their
component augmentations, which maps to zero because their target points
are joined by a path. Thus the map is not an $E$-equivalence.
The program verifies the permutation products, orbit sets and exponent-sum
matrix. The infinite group itself is specified by reduced words with free
reduction as multiplication, so no finite quotient is being substituted for
the stated category. This confirms the inference proposed after PATHS.
\end{proof}

\begin{claim}{B.12}{Proved}{A finite bifibration defeats both $W_1$ and rational local-system equivalences}
Put $W_{\mathbb Q}^{\mathrm{lc}}=W_{D(\mathbb Q)}^{\mathrm{lc}}$, where
locally constant diagrams are coherent diagrams of unbounded rational
complexes. There is a finite bifibration $p:H\to K$, with $14$ objects and
$100$ arrows upstairs, which belongs to both
$\Liss_W\cap\Prop_W$ for every fundamental localizer $W$, but belongs to
neither $\Fib_{W_1}$ nor $\Fib_{W_{\mathbb Q}^{\mathrm{lc}}}$.
The functor $p$ itself belongs to $W_{\mathbb Q}^{\mathrm{lc}}$.
\end{claim}
\begin{proof}
Let $K$ be the poset of nonempty proper subsets of $\{0,1,2,3\}$, with
inclusions as arrows. For composable arrows $u:x\subseteq y$, $v:y\subseteq z$,
put $c(u,v)=1\in C_2$ exactly for
\[
 (x,y,z)=(\{0\},\{0,1\},\{0,1,2\}),
\]
and put $c(u,v)=0$ otherwise. A triple containing an identity satisfies the
normalized cocycle equation by cancellation; there is no triple of strict
inclusions. Hence $c$ is a normalized $2$-cocycle. Define $H$ to have the
objects of $K$, arrows $(u,m)$ for $m\in C_2$, and composition
\[
 (v,n)(u,m)=(vu,m+n+c(u,v)),\qquad 1_x=(1_x,0).
\]
These formulas specify every identity and composite; the cocycle equation
proves associativity. Projection to $u$ defines $p:H\to K$.
For a prescribed lift of a base arrow, every cartesian or cocartesian
factorization is obtained by solving a unique equation in $C_2$.
Thus $p$ is a bifibration, with fibre $BC_2$ and identity transports on
its automorphism group. Its lifting categories have terminal or initial
objects, proving the two assertions about smoothness and properness for
the same arbitrary $W$.

Here is a nonabelian fundamental-group calculation. Write
$t_x=(1_x,1)$ and $h_{xy}=(x\subseteq y,0)$. Groupoid completion imposes
\[
 t_x^2=1,\qquad t_yh_{xy}=h_{xy}t_x,\qquad
 h_{yz}h_{xy}=h_{xz}t_x^{c(x,y,z)}.
\]
After collapsing a spanning tree, all the $t_x$ become a single central
involution $t$. The nerve of $K$ is the triangulated $2$-sphere.
Removing the interior of the one triangle supporting $c$ leaves a disk
with the same one-skeleton. Its triangle relations kill all remaining
horizontal generators. The omitted triangle relation then reads $t=1$.
Consequently $BK$ and $BH$ are connected and simply connected.
For the strict pullback $H\times_KH$, replace $C_2$ by $C_2^2$ and $c$
by $(c,c)$. The identical calculation leaves
$C_2^2/\langle(1,1)\rangle\cong C_2$.
The finite program checks the $1{,}552$ associative triples of $H$, the
$12{,}416$ triples of this pullback, all cartesian and cocartesian
factorizations, and performs Tietze eliminations in the full group
presentations. Its outputs are respectively the trivial group and
$\langle t\mid t^2=1\rangle$; it does not substitute first homology
for a fundamental-group computation.

For $W_1$, pull back the $W_1$-map $e\to K$ along $p$. The resulting
map $BC_2\to H$ fails to induce an isomorphism on fundamental groups.
This already proves $p\notin\Fib_{W_1}$.
For the rational assertion we specify the coefficient criterion being used.
If $Z$ is connected and based, evaluation at the basepoint on
$\Fun(Z,D(\mathbb Q))$ is conservative and preserves colimits. Its left
adjoint has monad $M\mapsto C_*(\Omega Z;\mathbb Q)\otimes_{\mathbb Q}M$.
The usual bar resolution therefore identifies the diagram category with
modules over this differential graded loop algebra. For a based map
$f:Z\to Z'$, restriction becomes restriction of scalars along
$C_*(\Omega f;\mathbb Q)$. It is an equivalence precisely when that
algebra map is a quasi-isomorphism: necessity follows by applying the
unit of extension--restriction to the regular module; sufficiency follows
by its bar resolution. The criterion is unchanged on homotopy categories,
since shifts detect all homotopy groups of mapping spectra.
By \citet[Theorem 3]{BCRiveraWierstraZeinalian2021}, this is equivalent
to an isomorphism on $\pi_1$ and rational isomorphisms on all higher
homotopy groups. Applied componentwise, it also requires a bijection on
components: a missing component prevents conservativity, and two components
mapping into one prevent essential surjectivity by choosing zero and
nonzero values on the two components.

By B.7 the strict fibre $BC_2$ is the homotopy fibre of $BH\to BK$.
Its long exact sequence shows that $\pi_i(p)$ is an isomorphism for
$i\geq3$ and has kernel and cokernel of exponent two in degree two.
Both fundamental groups vanish. The preceding rational criterion proves
$p\in W_{\mathbb Q}^{\mathrm{lc}}$.
Pulling this very map back along itself gives
$H\times_KH\to H$, whose fundamental groups are $C_2\to1$.
It is not in $W_{\mathbb Q}^{\mathrm{lc}}$; explicitly the sign local
system on the source cannot be restricted from the simply connected target.
This proves the second intrinsic-fibration failure with finite categories.
\end{proof}

\begin{claim}{B.13}{Proved}{The covering witnesses have different behaviour for local coefficients}
Every covering of ordinary categories, meaning a discrete fibration and a
discrete opfibration, belongs to $\Fib_{W_1}$ and to
$\Fib_{W_{\mathbb Q}^{\mathrm{lc}}}$. In particular the three-sheeted
covering of B.11 belongs to these two intrinsic classes, although it
belongs to no $\Fib_{W_E}$ for $E\ne0$.
The endpoint map of B.8 with $C=BC_2$ lies in the same two intrinsic
classes, while it fails to lie in the rational-homology intrinsic class.
\end{claim}
\begin{proof}
A covering has discrete strict fibres and bijective transports.
Its nerve realizes a homotopy covering, and after any strict base change its
strict pullback is the homotopy pullback, by B.7. On each connected component
a covering is classified by a $\pi_1$-set. A $W_1$-map is bijective on
components and induces isomorphisms of fundamental groups, so its pullback
identifies orbit sets and corresponding stabilizer groups. It therefore
induces a $W_1$-map of the total coverings. For
$W_{\mathbb Q}^{\mathrm{lc}}$, use the criterion proved in B.12 in addition:
a covering induces isomorphisms on all homotopy groups in degrees at least
two, and hence the corresponding total-space map is a rational isomorphism
in those degrees. This proves universal preservation, including disconnected
or empty fibres and coverings with infinitely many sheets.

The map $E_x(BC_2)\to BC_2$ has contractible total space and homotopy
fibre the discrete group $C_2$, so it represents the same homotopy covering
as the universal two-sheeted cover. All its strict pullback squares compute
homotopy pullbacks by B.7. The preceding argument therefore applies to it,
without claiming that the actual categorical functor is a covering.
The input $e\to BC_2$ is a rational-homology equivalence, since averaging
contracts positive-degree rational group homology. Its pullback along the
endpoint map is the inclusion of a fibre with two contractible components
into a contractible total category; on $H_0(-;\mathbb Q)$ it is
$\mathbb Q^2\to\mathbb Q$. This proves the failure for rational homology.
In contrast, $e\to BC_2$ is neither a $W_1$-map nor a
$W_{\mathbb Q}^{\mathrm{lc}}$-map. Likewise the free-group endomorphism
of B.11 is not a $W_1$-map: it changes the orbit set of the specified action.
Restriction of the corresponding rational permutation representation changes
its invariant-vector dimension from one to two, so it is not a
$W_{\mathbb Q}^{\mathrm{lc}}$-map either. The coefficient distinction is
therefore essential to each use of the example.
\end{proof}

\begin{claim}{B.14}{Refuted}{Cohomological fibrations are not the proper class}
Let $\mathcal D$ be the derivator of a nonterminal presentable
infinity-category, and use presheaf variance. Define a cohomological
$\mathcal D$-fibration by requiring, after every strict base change, that
right Kan extension sends every locally constant coefficient to a locally
constant coefficient and has the canonical strict-fibre formula.
This class and $\Prop_{W_{\mathcal D}^{\mathrm{lc}}}$ are not contained
in one another. Dually, the homological class and
$\Liss_{W_{\mathcal D}^{\mathrm{lc}}}$ are not contained in one another.
In particular both conclusions apply to sets, spaces, and rational complexes.
\end{claim}
\begin{proof}
Use first the map $p:V\to[1]$ in B.7. Its section $s$ is left adjoint to
$p$: an arrow $s(t)\to v$ is uniquely the arrow of $[1]$ from $t$ to $pv$.
This adjunction has identity unit, and its counit $sp\Rightarrow1$ lies
over the base. The same formulas define an adjunction after every strict
base change $p_T:V_T\to T$. If $L$ is locally constant on $V_T$, the
counit induces an equivalence $L\simeq p_T^*s_T^*L$, since its components
are arrows of $V_T$. Put $F=s_T^*L$, a locally constant object on $T$.
The unit $F\to(p_T)_*p_T^*F$ is an equivalence. Indeed, at $t\in T$
its comma formula is induced by
\[
 V_T/t\longrightarrow T/t.
\]
This functor has the left adjoint induced by $s_T$, so its opposite is
initial for limits; the resulting limit is $F_t$. Equivalently, apply
the adjunction contraction to the bar resolution of that limit.
Consequently $(p_T)_*L\simeq F$ is locally constant. The strict fibres
are $\{a<b\}$ or $\{c\}$, both contractible. Restricting the same
unit to a fibre identifies its limit with $F_t$. These are the canonical
comparison maps, natural under any further base change, which proves the
required universal property.
The outgoing lifting category from $b$ over $0\to1$ is empty, so $p$
is not $W_{\mathcal D}^{\mathrm{lc}}$-proper. The empty category is not
$W_{\mathcal D}^{\mathrm{lc}}$-aspherical: restriction from the
nonterminal category $\mathcal D(e)$ to the terminal category of diagrams
on the empty category is not an equivalence.

Conversely, the cosieve $j:\{1\}\to[1]$ is an opfibration and therefore
$W$-proper for every $W$. Choose a nonterminal coefficient $M$. Its right
Kan extension is the presheaf with values the terminal object at $0$ and
$M$ at $1$; the transition $M\to1$ is not an equivalence. Thus even this
constant input has non-locally-constant right direct image. This proves the
other failed containment. Reversing all indexing categories and exchanging
limits with colimits gives the homological assertion, using a noninitial
coefficient. A nonterminal presentable category has both a nonterminal and
a noninitial object.
This concerns the universal condition on all local coefficients, not merely
the direct image of one constant coefficient. The two tests must not be
interchanged in the manuscript's Chapter~I arguments.
\end{proof}

\begin{claim}{B.15}{Proved}{Disposition of the inherited finite coefficient examples}
The following table evaluates the maps specified in GALOIS D.3--D.5 for
$W=W_1$ and $W=W_{\mathbb Q}^{\mathrm{lc}}$. A pair of entries gives
these two localizers in that order; a single entry applies to both.
Here ``smooth'' and ``proper'' always use the same $W$ as the intrinsic
column.
\[
\begin{array}{c|c|c|c}
\text{specified map}&W\text{-smooth}&W\text{-proper}&\Fib_W\\\hline
\mathrm{D.3}:p:E\to B&\text{yes}&\text{yes}&\text{yes}\\
\mathrm{D.3}:q:B\to[1]&\text{no}&\text{yes}&\text{no}\\
\mathrm{D.3}:qp&\text{no}&\text{yes}&\text{no}\\
\mathrm{D.4}:p:E\to B&\text{yes}&\text{yes}&\text{no}\\
\mathrm{D.4}:q:B\to[1]&\text{no}&\text{yes}&(\text{yes},\text{no})\\
\mathrm{D.4}:qp&\text{no}&\text{yes}&\text{no}\\
\mathrm{D.5}:H\to K&\text{yes}&\text{yes}&\text{no}\\
\mathrm{D.5}:q:R\to[1]&\text{no}&\text{yes}&\text{no}
\end{array}
\]
\end{claim}
\begin{proof}
We give the constructions, so the table does not use the earlier answer as
an authority. In D.3, $B$ has objects $a,b$, $B(a,a)=\{1\}$,
$B(b,b)=C_2$, $B(a,b)=C_2$, and $B(b,a)=\varnothing$; postcomposition
on $B(a,b)$ is addition in $C_2$. The projection $q$ is the opfibration
of $e\to BC_2$ over $[1]$. Put $L(a)=L(b)=C_2$, with every labelled
arrow acting by addition, and $E=\int_B L$. Thus $p$ is a covering and
B.13 applies. For $q$, the incoming lift category of $0\to1$ at $b$ is
a discrete two-point category. For $qp$ the corresponding category is
again discrete with two objects. Neither is aspherical for either localizer.
All three maps are opfibrations and hence proper. The inclusion of the
fibre at $0$ into $B\simeq BC_2$ fails on $\pi_1$, and the inclusion
of the fibre $C_2^{\mathrm{disc}}$ into $E\simeq EC_2$ fails on
components. Pulling back $\{0\}\to[1]$ gives exactly these witnesses
for the two failures of intrinsic membership.

For D.4, take $K,H$ from B.12 and $x_0=\{0\}$. The category $B$ has
$K$ and a new object $a$, arrows $B(a,x)=K(x_0,x)$, no arrows into $a$
except its identity, and the compositions forced by $K$. Let
$r:B\to K$ send $a$ to $x_0$, let $q$ be the opfibration of
$e\xrightarrow{x_0}K$, and put $E=B\times_KH$. The map $p:E\to B$
is a bifibration. Its pullback along $K\hookrightarrow B$ is $H\to K$,
so B.12 excludes intrinsic membership. The retraction $r$ and the arrows
$a\to x_0$ give $BB\simeq BK\simeq S^2$; their label-zero lifts give
$BE\simeq BH$. The fibre of $qp$ at $0$ is $BC_2$, so its inclusion in
$E$ fails both localizers on fundamental groups. Objects $x$ of $K$ not
containing $0$ have no incoming lift from colour $0$, excluding smoothness
of $q$ and $qp$. Both are opfibrations.
For every $T\to[1]$, the fibres of $q_T$ are a point or the simply connected
category $K$. In the groupoid completion of its Grothendieck construction,
every fibre object has a unique path from the chosen section object.
Factorization into a cocartesian arrow followed by a fibre arrow proves
naturality of these paths. Thus $\Pi_1(q_T)$ is an equivalence, for every
$T$, and two-out-of-three gives $q\in\Fib_{W_1}$.
For rational local systems, the pullback of $\{0\}\to[1]$ is
$e\to B\simeq S^2$, which fails on $\pi_2\otimes\mathbb Q$.

Finally $R$ has objects $b,a$, automorphism group $C_2$ at $b$, identity
only at $a$, one arrow $b\to a$, and the relation that precomposition
by either automorphism fixes that arrow. Its projection is the opfibration
of $BC_2\to e$. The incoming lift category at $a$ is $BC_2$, excluding
smoothness for both localizers. The category $R$ has terminal object $a$;
the fibre inclusion $BC_2\to R$ fails on $\pi_1$ and detects the
intrinsic failure. The first D.5 map is precisely B.12.
The code specifies these finite categories by tables and checks their units,
associativity, all relevant lift categories and component data. Thus the
proofs retain the different roles of weak constant, fully faithful, and
equivalent local-coefficient tests.
\end{proof}

\begin{claim}{B.16}{Proved}{Accessible coefficient localizations and the limit of the general statement}
If $W$ is an accessible fundamental localizer, let $\Spaces_W$ be the
presentable localization of spaces corresponding to $W$. Then
\[
 (\Cat/S)[(W^{\mathrm{pre}}_S)^{-1}]_\infty
 \simeq\Fun(NS^{\op},\Spaces_W),\qquad
 (\Cat/S)[(W^{\mathrm{co}}_S)^{-1}]_\infty
 \simeq\Fun(NS,\Spaces_W).
\]
For an arbitrary fundamental localizer the relative definitions, the
opposition isomorphism, B.6 and B.9 remain valid without accessibility.
A presentable coefficient category in the displayed formulas is not
asserted without the accessibility hypothesis.
\end{claim}
\begin{proof}
For a strict diagram $F:S\to\Cat$, write $\int_SF$ for the opfibred
Grothendieck construction. The comma functor $R_S$ of B.6 and this
construction are relative functors between
$(\Cat/S,W^{\mathrm{co}}_S)$ and the category of strict $S$-diagrams
with pointwise $W$-maps. Indeed, $R_S$ detects its specified class by
its definition, while $(\int_SF)/s\to F(s)$, given by cocartesian
transport, is left adjoint to the fibre inclusion. It is consequently in
$W$, naturally in $F$ and $s$. The natural unit
$X\to\int_SR_SX$ is fine by B.6, and the natural evaluation
$R_S\int_SF\to F$ is pointwise in $W$. Their localizations are inverse.
For accessible $W$, the combinatorial Thomason model category of
\citet[Theorem 5.2.15]{BCCisinski2006} presents $\Spaces_W$.
Its category of strict diagrams has the derived diagram comparison:
\citet[Example 7.9.6 and Theorem 7.9.8]{BCCisinski2019} apply using
functorial model factorizations and small homotopy colimits. Thus its
pointwise localization is $\Fun(NS,\Spaces_W)$. This verifies the
strict-diagram step, rather than inferring it from an absolute equivalence
of homotopy categories. Opposition proves the first formula.
The hypotheses of the purely relative arguments in B.1, B.6 and B.9 do
not involve the model structure or a small generating set, which proves
the last assertion. For example $W_0$ gives set-valued diagrams, whose
coarse localization need not identify with their locally constant part;
Part~E treats that separate obstruction.
\end{proof}

\begin{claim}{B.17}{Known}{Intrinsic fibrations for truncation localizers}
For $n\geq0$, a functor $p:T\to B$ lies in $\Fib_{W_n}$ if and only if,
for every $X\to B$, the canonical map
\[
 \tau_{\leq n}B(T\times_BX)\longrightarrow
 \tau_{\leq n}BT\times_{\tau_{\leq n}BB}\tau_{\leq n}BX
\]
is an equivalence of $n$-types. The pullback on the right is a homotopy
pullback in $\Spaces_{\leq n}$.
For $n=0$ this is equivalent to the following elementary pair of conditions.
For every connected component $B_j$ of $B$, write $T_j=T\times_BB_j$.
Then $\pi_0(T_b)\to\pi_0(T_j)$ must be bijective for each $b\in B_j$;
and for every arrow $\beta:b\to b'$ and every component $c$ of $T_j$
there must be an arrow of $T$ over $\beta$ in component $c$.
\end{claim}
\begin{proof}
The general equivalence is precisely the intrinsic weak-fibration criterion
of \citet[Proposition 6.2.3]{BCCisinski2006}. Its proper-localizer
hypothesis is supplied for $W_n$ by that book's \S9.2.1, and
Corollary~9.2.15 identifies its weak equivalences with the asserted
truncations. These hypotheses matter: the criterion is not being applied
to a nonproper homology localizer.
For clarity, the elementary specialization can also be proved directly.
Pullback of $e_b\to B_j$, a $W_0$-map, gives the first condition.
Pullback of either endpoint of $[1]\xrightarrow{\beta}B$ then forces
an arrow joining the two layers in each prescribed component, giving the
second. Conversely, for connected $X_a\subset X$, label objects of
$T\times_BX_a$ by their component $c$ in $T_j$. Objects over any fixed
$x\in X_a$ with a given label form a nonempty connected category by the
first condition. Every arrow of $X_a$ has a lift connecting the associated
labelled categories by the second. Finite zigzags therefore connect exactly
the objects with the same label. Hence, naturally in $X$,
\[
 \pi_0(T\times_BX)\cong\pi_0(T)\times_{\pi_0(B)}\pi_0(X).
\]
This formula preserves $W_0$-maps and proves sufficiency, including empty
fibres. It also explains directly why the regular $C_2$ covering fails B.10.
\end{proof}

\begin{claim}{B.18}{Proved}{Useful bounds for a homology intrinsic class}
For every spectrum $E$, $\Fib_{W_E}$ contains product projections and
inclusions of unions of connected components, together with their closure
under composition, strict base change and retracts. It is contained in the
class of $p:T\to B$ for which every inclusion
\[
 T_{\sigma(i)}\longrightarrow T\times_B[m],
 \qquad \sigma:[m]\to B,quad 0\leq i\leq m,
\]
is an $E$-equivalence. For $E\ne0$ this upper bound is strict.
\end{claim}
\begin{proof}
For a product projection the induced map is $u\times1_K$.
The identity $(X\times Y)_+\cong X_+\wedge Y_+$ identifies its
$E$-homology comparison with the smash product of the $E$-equivalence
of $u$ and $\Sigma^\infty(BK)_+$. This is an equivalence.
For a component inclusion the nerve is a disjoint union of components;
restriction selects summands of a componentwise equivalence.
The three closure assertions follow from iterated pullbacks and retracts.
The inclusion $\{i\}\to[m]$ is a weak homotopy equivalence, so its
pullback gives each necessary test in the display.
The covering of B.11 passes every one of these tests: over a simplex it
is a constant covering, each fibre inclusion being a weak homotopy
equivalence. B.11 proves that it nevertheless fails intrinsic membership
for every nonzero $E$. This proves strictness without identifying model
fibrations of a nonproper localization with intrinsic fibrations.
\end{proof}

\begin{claim}{B.19}{Proved}{The full fence construction and its mapping-path factorization}
Let $\mathcal P(C)$ be the category of all eventually constant diagrams
$J\to C$ on the integer fence of B.8, with all natural transformations.
Write $\sigma,\beta$ for the eventual left and right values. Then
$(\sigma,\beta):\mathcal P(C)\to C\times C$ is smooth and proper
for $W_\infty$, and therefore for every fundamental localizer.
For every $f:A\to C$, the functors
\[
 E(f)=A\times_{C,\sigma}\mathcal P(C),\qquad
 i_f(a)=(a,\operatorname{const}_{f(a)}),\qquad
 p_f(a,\gamma)=\beta(\gamma)
\]
give $f=p_fi_f$, with $i_f\in W_\infty$ and $p_f$ smooth and proper.
All formulas are natural in strictly commuting squares.
For $C=[1]$, the endpoint functor is neither a Grothendieck fibration nor
an opfibration. Thus the verified fence construction in PATHS has exactly
the asserted homotopical properties, without being a bifibration.
\end{claim}
\begin{proof}
For an outgoing endpoint pair $(\alpha,\beta)$ from a path $\gamma$,
choose a sufficiently negative odd integer $a$ and sufficiently positive
even integer $b$. Replace the left tail through $a$ by the source-target
of $\alpha$ and join it at $a+1\to a$ using $\alpha$; replace the
right tail after $b$ by the target of $\beta$ and join it at $b\to b+1$
using $\beta$. There is a natural transformation from $\gamma$ to this
new path, identity between the joins and the prescribed endpoint arrows
outside. Among outgoing lifts stationary outside this interval, it is
initial: the comparison to any other lift is that lift's transformation
inside, and the endpoint identities outside. Naturality at the two joins
is precisely the two original commutative squares. Uniqueness follows from
the identity endpoint components and stationary tails.
Letting the interval expand gives a filtered union of these contractible
categories; every simplex and every finite homotopy in the nerve belongs
to a stage. Thus the entire lifting category is weakly contractible.
For incoming lifts use an even left boundary and an odd right boundary;
the same argument gives terminal stage objects. This proves smoothness
and properness by the exact lifting tests. It is the two-ended version
of B.8, with the parity choices specified.

The functor $\sigma$ is the composite of this endpoint map and a product
projection, so is also smooth and proper. Its strict fibre at $x$, with
identity left endpoint on morphisms, is exactly $E_x(C)$ of B.8 and is
weakly contractible. B.7 now shows that every strict base change of
$\sigma$ is a weak equivalence. In particular $E(f)\to A$ is one;
its section $i_f$ is a weak equivalence by two-out-of-three.
The map $E(f)\to A\times C$ is a base change of the full endpoint map,
and projection $A\times C\to C$ is both smooth and proper. Stability
under composition proves the assertion for $p_f$. B.7 also identifies
its strict fibres with the homotopy fibres of $Bf$. Applying a commuting
square pointwise to these diagrams proves strict naturality.

For the failure of a fibration, incoming lifts into the constant $1$ path
over $(0\to1,1)$ are the eventual $[1]$-paths from $0$ to $1$, ordered
pointwise. A proposed terminal path is still $0$ sufficiently far left;
a single-transition path changing to $1$ earlier has no map to it.
Dually, in the outgoing lift category from constant $0$ over
$(1,0\to1)$, a proposed initial path changes to $1$ before some
later single-transition path and has no map to that later path.
Single transitions are taken across an even-to-odd forward edge, so all
these diagrams are legitimate fence functors. The obstructions hold at
arbitrarily large supports, which is why a finite cut-off cannot establish
existence of universal lifts. This proof verifies the exact convention of
PATHS E.4--E.5; it does not identify it with every convention called an
Evrard cocylinder \citep{BCEvrard1975}.
\end{proof}

\begin{claim}{B.20}{Refuted}{Reversing path words does not repair the all-type endpoint functor}
Use the all-type category $\operatorname{Ch}(C)$ with objects finite
oriented interval diagrams and morphisms $(h,\eta)$, where $h$ is an
endpoint-preserving, weakly monotone functor of their interval types and
$\eta:\gamma\Rightarrow\delta h$. Its endpoint map is smooth for every $C$, but for $C=[1]$
it is not proper and is neither a Grothendieck opfibration nor a fibration.
The conclusion cannot be transferred from B.19 by calling both objects
``path categories'' or by reversing the path word.
\end{claim}
\begin{proof}
First fix a path $\gamma$ of type $\tau$ and incoming endpoint arrows
$\alpha:a'\to a$, $\beta:b'\to b$. Adjoin these arrows to obtain the
path $\delta$ of type $(+,\tau,-)$ from $a'$ to $b'$, with its collapse
map $d:\delta\to\gamma$. Let $L$ be the incoming lifting category
and let $R=\operatorname{Ch}(C;a',b')/\delta$. Composition with $d$
defines $i:R\to L$, and $R$ has terminal object $1_\delta$.
For a lift $(h,\eta):\zeta\to\gamma$, pad $\zeta$ with an identity
edge at each end. The extended type map $(+,h,-)$ and the transformation
whose middle is $\eta$ and whose new endpoints are identities define
an arrow of this padded path to $\delta$. The joining naturality equations
are precisely $\eta_{\mathrm{left}}=\alpha$ and
$\eta_{\mathrm{right}}=\beta$. This defines a functor $F:L\to R$.
Collapsing the two added identity edges gives a natural transformation
$iF\Rightarrow1_L$ over $\gamma$. Hence $BL$ is homotopy dominated
by the contractible space $BR$, and is itself contractible. This proves
smoothness, with the exact type convention retained.

Let $e$ be the empty interval type. An endpoint-preserving type map
$e\to\tau$ forces the two endpoints of $\tau$ to coincide, so $\tau=e$.
For the empty path at $0$, the base arrow
$(1_0,0\to1):(0,0)\to(0,1)$ consequently has no outgoing lift.
Its outgoing lift category is empty. This proves failure of properness
and of an opfibration, for every nontrivial component-detecting localizer.
For failure of a fibration use the empty path at $1$ and the incoming
base arrow $(0\to1,1_1)$. Its incoming lift category is
$\operatorname{Ch}([1];0,1)$, since collapse to $e$ and the transformation
to constant $1$ are unique. Suppose a terminal path $c$ had nonempty
type $\tau$. Form $d$ on the concatenated type $\tau\tau$, constant
at $0$ on the first block and equal to $c$ on the second.
There are two distinct type maps $\tau\tau\to\tau$, collapsing
respectively the first and the second block. The first makes $d$ exactly
the pullback of $c$. For the second pullback, the first block is $c$ and
the second constant at $1$; the pointwise inequalities define a natural
transformation from $d$. Both transformations have identity endpoints.
Hence there are two distinct arrows $d\to c$, contradicting terminality.
All identities and composites in this category are respectively
$(1,1)$ and $(kh,(\theta h)\eta)$, so the argument specifies the
actual infinite counterexample and its type-map distinction.
Taking the opposite of a Grothendieck construction reverses its indexing
category of types. Reversing each interval word within the unchanged
indexing category does not perform that operation. This locates the exact
variance error in the proposed deduction of properness at VII ms.~161--166.
The two distinct constructions therefore confirm the classification in
PATHS E.3--E.5. The factorization used here is the independently verified
fence factorization of B.19.
\end{proof}

\begin{claim}{B.21}{Open}{Replacing a smooth and proper representative by a bifibration}
For the one-object category of the monoid
$R=\{1,e_{00},e_{01},e_{10},e_{11}\}$ with
$e_{ij}e_{kl}=e_{il}$ and $1r=r1=r$, it remains unsettled here whether
the local system of path spaces has a representative by a strict
Grothendieck bifibration with contractible total category.
The eventual fence endpoint functor is not such a representative.
\end{claim}
\begin{proof}[Exact remaining issue and verified obstruction]
Put $e=e_{00}$. Both $eR$ and $Re$ consist of nonunits.
For incoming lifts to a constant fence path over endpoint pair $(1,e)$,
a natural transformation has eventual components $1$ on the left and
$e$ on the right. A candidate terminal lift has some component equal to
$e$. Move the single change from $1$ to $e$ farther right, across an
odd-boundary incoming edge, to obtain another lift whose component there
is still $1$. A map to the candidate would require $e\theta=1$, impossible.
For outgoing lifts the same construction across an even boundary would
require $\theta e=1$, again impossible. These equations verify the
failure in the precise monoid, whose $125$ associativity equations and
nonunit ideals are checked by the program.
They do not obstruct a different total category equipped with coherent
cartesian and cocartesian universal arrows. Constructing that representative,
or an invariant forbidding every such representative, is the remaining
statement. No bifibration-existence assertion is used in Parts~B or~E.
\end{proof}

\begin{claim}{B.22}{Proved}{The maximal simplicial classes and the all-base recipe}
For every simplicial base $S$, let $W^{\mathrm{cov}}_S$ and
$W^{\mathrm{contra}}_S$ denote the actual covariant and contravariant
equivalences. The maps whose every base change preserves these classes
under raw pullback are, respectively, Cisinski-smooth and Cisinski-proper
maps. Moreover the manuscript's universal test recipe gives exactly
\[
 W^{\Liss_{\mathrm{Cis}}}_S=W^{\mathrm{cov}}_S,
 \qquad W^{\Prop_{\mathrm{Cis}}}_S=W^{\mathrm{contra}}_S.
\]
Thus opposition gives a perfect duality over $S^{\op}$ for every simplicial
base, although B.4 refutes the version using literal slices at an arbitrary
base. For accessible $W$, replace fibrewise equivalences of spaces by
fibrewise $W$-equivalences after left or right fibrant replacement.
These two accessible localizations classify diagrams with values in
$\Spaces_W$ on the directed infinity-category presented by $S$ or its
opposite, and their restrictions to ordinary nerves are B.16.
Their universal raw-pullback classes contain the minimal smooth and proper
classes. Complementarity with intrinsic total $W$-equivalences holds for
these classes only for $W_\infty$ and the total localizer.
\end{claim}
\begin{proof}
Pullback along a Cisinski-smooth map is left Quillen for the covariant model
structures \citep[Proposition 4.4.18]{BCCisinski2019}; this remains true
after every base change. Conversely, suppose a simplicial map $p$ has the
stated universal preservation property. If $i:X\hookrightarrow Y$ is
left anodyne and $Y$ maps to the base of $p$, regard $i$ as a covariant
weak equivalence over its own target $Y$. The pullback $i'$ is a monomorphism
and a covariant weak equivalence over its own target. Such a map is left
anodyne: the trivial-cofibration lifting property in that slice applies to
every left fibration over the target, which is precisely the defining left
anodyne lifting property. Consequently all pullbacks of left-anodyne maps
along $p$ are left anodyne. The opposite of
\citet[Proposition 4.4.4]{BCCisinski2019} identifies this condition with
Cisinski-smoothness. This proves maximality.
A covariant equivalence passes every smooth pullback test and then is a
total weak equivalence. Conversely the tests include the right-fibrant
replacements of each vertex, which are smooth right fibrations.
B.2's all-base formulation detects covariant equivalences on precisely
those tests. Opposition proves the entire contravariant statement.

For accessible $W$, compose the covariant classifier with pointwise
localization $\Spaces\to\Spaces_W$. This is an accessible reflective
localization, computed pointwise, and its inverse image specifies the
claimed weak equivalences. A choice of a categorical fibrant replacement
of $S$ identifies its target with diagrams on the infinity-category
presented by $S$; no literal $\Fun(S,\Spaces_W)$ with a nonfibrant
simplicial set as infinity-category is used.
If $p$ is minimally smooth, replace an arbitrary $X\to S$ by a left
fibration $X^{\mathrm{lf}}\to S$. The unit remains a minimal covariant
equivalence after pullback along $p$. For left fibrations the new weak
equivalences are fibrewise $W$-equivalences, and arbitrary pullback preserves
these tests exactly. Two-out-of-three proves preservation for arbitrary
$X,Y$. This proves containment of the minimal smooth class, and its dual.
If their intersection universally preserved total $W$-equivalences, the
nerves of the minimal smooth and proper factorizations in B.9 would imply
homotopy-pullback stability of $W$. The proof of B.9 then forces $W_\infty$
or the total class. The two converse cases follow from B.7 or vacuity.
In particular every finite and homological counterexample above gives a
simplicial counterexample simply by taking nerves. This argument does not
rename a general $W$-smooth map ``Lurie-smooth'': Lurie's unqualified term
uses the minimal space-valued theory.
\end{proof}
